% !TeX root = ../../quotients-and-descent.tex
\section{Introduction}
This paper asks an integral question within one complex CM isogeny class:
which quotients of a fixed sixfold admit principal polarizations, and over
which fields can the quotients and their polarizations be defined? A
rational Hodge-algebraicity theorem does not answer either question.
The starting variety is defined from a curve, independently of any
Hilbert-modular packet.

Let $H$ be the smooth projective model of $y^2=x^{21}-1$. Its primitive
$\Phi_{21}$ factor $A\subset J(H)$ is six-dimensional and has full
endomorphism ring $\mathcal O_M$, where $M=\mathbb Q(\zeta_{21})$.
We keep this actual embedded factor distinct from isogenous principal
models. Its restricted Jacobian polarization has degree $35721$, whereas
its minimum polarization has degree $729$ and type $(1,1,1,3,3,3)$.

\subsection{The finite family and the equivalence relations}
Fix the minimum polarization $\lambda_{\min}$ and write its kernel as
$A[\pi]\simeq\mathbb F_3^6$, with the notation made explicit below.
The family consists of $A$ and all quotients $A/H_0$ for subgroups
$H_0\subset A[\pi]$ of order $27$. All isomorphisms preserve the origin.
An underlying class forgets the polarization; a polarized class retains
it but forgets CM and RM markings. A kernel is a subgroup of the fixed
source, not a target isomorphism class. These distinctions account for the
different counts in the following summary.

\begin{center}
\begin{tabular}{lr}
\toprule
Object being counted & Number\\\midrule
Order-$27$ subgroups of the fixed source kernel & $33880$\\
Underlying quotient classes (excluding $A$) & $188$\\
Underlying classes including $A$ & $189$\\
Principalizable underlying quotients & $80$\\
Principal polarized quotient classes & $160$\\
Nonprincipalizable underlying quotients & $108$\\
Minimum polarized classes including $A$ & $2969$\\\bottomrule
\end{tabular}
\end{center}
The last number is $1+160+108\cdot26$. The source contributes its unique
minimum polarization, and each obstructed quotient contributes $26$.
Fixing the pullback form gives $1120$ principal-target kernels; allowing
all minimum pullback forms gives $14224$. Neither number counts all
$33880$ order-$27$ subgroups as principal-target kernels.

\subsection{Main results and proof route}
Theorem~\ref{pp:thm:fermatallquotients} classifies the principal targets;
Theorem~\ref{pp:thm:fermatunderlying} forgets the polarizations and identifies
the obstruction; Theorem~\ref{pp:thm:obstructedminimum} determines the
minimum forms on the remaining $108$ varieties. Their proofs pass from
integral lattices to finite residue subspaces and distinguish the relevant
unit actions before taking orbits.

The Galois action then determines fields of moduli, but a fixed isomorphism
class need not descend. Put
\[
 \mathcal E=M\bigl(\sqrt{(3+\sqrt{21})/2}\bigr).
\]
Theorem~\ref{pp:thm:quotientdefinitionfield} proves that for number fields $M\subseteq F\subset\overline{\mathbb Q}$, any one of the $188$ quotient classes has an $F$-model
exactly when $F$ contains $\mathcal E$. A single twist makes all subgroup
kernels and quotient maps available over $\mathcal E$. Thus the common
field-of-moduli computation and the actual defining-field construction
are distinct stages of the argument.
Theorem~\ref{pp:thm:definitiondegrees} gives sharp absolute degrees for
$168$ quotient classes and lower bounds for the other $20$.

The final geometric section treats Jacobians. It distinguishes the
restricted ten-class principal census with $R$-containment from the $156$
principal classes with smaller endomorphism order $S$. The former has
exactly one canonical Jacobian under the stated enclosure and execution
contract. The latter $156$ statuses remain open.

\subsection{Arithmetic inputs and relation to other work}
Write $z=\zeta_{42}=-\zeta_{21}^{11}$, $t=z+z^{-1}$, $L=M^+$ and
\[
 f(T)=T^6+T^5-6T^4-6T^3+8T^2+8T+1,\qquad
 \mathfrak P=(1-z^{14})\mathcal O_M,
 \quad R=\mathcal O_L+\mathfrak P,\quad S=\mathbb Z+\mathfrak P.
\]
Here $\mathfrak P$ is the prime over $3$; its residue field is
$\mathbb F_{729}$ and $R/\mathfrak P=\mathbb F_{27}$. The letter $\pi$
below is $1-z^{14}$, a generator of $\mathfrak P$ in $M$. It is \emph{not} the norm-$43$ generator of
Paper I. We use $K=\mathbb Q(\sqrt{-3})$ for the Weil action.
Write $E_L=\mathcal O_L^\times$ and let $E_L^+$ denote its subgroup
of totally positive units.

The imported arithmetic inputs are the ideal classification and marked
polarization classification of Paper I, Propositions~\externalref{packet}{pp:prop:picardorder},
\externalref{packet}{pp:prop:idealclassmonoid} and~\externalref{packet}{pp:prop:markedpolarizations},
and the unit-norm theorem in the supplement. They are established without
using this quotient family. The notations $B_0,B_1,C_0,C_1$ mean the two
principal-polarization classes on each of the two $R$-ideal models in the
B/C CM-type orbit; Paper I identifies them with its principal curve fibers.
The quotient construction here needs their arithmetic classification, not
that curve placement. The type representatives are
$\Phi_B=\{1,5,23,25,29,31\}$ and
$\Phi_C=\{1,5,11,13,17,19\}$; $\Phi_0=\Phi_C$.

Goodson and Saller discuss the rational CM-type/Fermat comparison
\cite{Goodson2024,Saller2026}. The general CM Hodge assertion in
\cite{OpenAICM2026}, if valid, concerns algebraicity of rational classes;
it does not itself provide the integral census or descent data here.
The existence of individual models must also be compared with classical
CM descent, including Milne's correction and addendum
\cite{MilneModuli1972,MilneModuli1974}. Our explicit simultaneous twist,
field obstruction and finite tables are distinguished from that existence
question. These comparisons are not a certification of historical priority.
